\section{Interpolation and Extrapolation Techniques}\label{sec:interpolation}

\subsection{The Runge Phenomenon in Spectroscopy (Polynomial vs. Rational/Spline Interpolation)}

In experimental spectroscopy, one often measures the intensity of a spectral line at discrete wavelength intervals. A typical absorption line shape is modeled by a Lorentzian (Cauchy) distribution:
$$ L(x) = \frac{1}{1 + 25x^2} $$
Consider the domain $x \in [-1, 1]$, representing a normalized frequency range around the line center.

\begin{enumerate}
    \item Sample this function at $N=11$ equidistant points (nodes) in the interval $[-1, 1]$. Perform a global polynomial interpolation of degree $N-1$ (Lagrange interpolation). Plot the original function and your interpolant on a fine grid.
    \item You should observe wild oscillations near the boundaries (the Runge phenomenon). Calculate the maximum error $\epsilon_{max} = \max |L(x) - P_{N-1}(x)|$.
    \item Repeat the interpolation using **Chebyshev nodes**, defined as $x_k = \cos\left(\frac{2k-1}{2N}\pi\right)$ for $k=1,\dots,N$. Plot the result and compare the maximum error.
    \item Finally, use a **Cubic Spline** on the original 11 equidistant nodes.
    \item \textit{Discussion:} If you were trying to measure the "equivalent width" (area under the curve) of this spectral line, which interpolation method would introduce significant unphysical artifacts, and why?
\end{enumerate}

\vspace{0.5cm}

\subsection{Thermodynamics of Stellar Interiors (Derivatives of Interpolants)}

In stellar structure simulations (e.g., White Dwarfs or Neutron Stars), the Equation of State (EOS) is often provided as a pre-computed table of Pressure ($P$) vs. Density ($\rho$). However, hydrodynamic stability calculations require the adiabatic sound speed, defined as:
$$ c_s^2 = \frac{dP}{d\rho} $$

\begin{enumerate}
    \item Generate a synthetic "EOS table" using the polytropic relation $P(\rho) = K \rho^{5/3}$ (modeling a non-relativistic degenerate Fermi gas). Create a dataset of 10 points logarithmically spaced between $\rho = 10^4$ and $\rho = 10^7$.
    \item \textbf{Linear Interpolation:} Interpolate $P(\rho)$ linearly between the points. numerically compute the derivative $dP/d\rho$ (the sound speed squared) from this interpolant on a fine grid. Plot $c_s^2$ vs $\rho$.
    \item \textbf{Cubic Spline:} Interpolate the same data using a Cubic Spline and compute the derivative of the spline. Plot this $c_s^2$ on the same graph.
    \item \textit{Discussion:} What physical violation occurs in the sound speed when using linear interpolation? Why is the Cubic Spline preferred for thermodynamic tables used in hydrodynamics codes?
\end{enumerate}

\vspace{0.5cm}

\subsection{Non-Parametric Inference with Gaussian Processes (Reconstructing the Universe)}

In modern data-driven physics (e.g., reconstructing the history of Dark Energy from Supernovae), we often need to interpolate functions where the data is noisy and we do not have a trusted physical theory to define the functional form. Gaussian Process (GP) Regression is the standard tool for this.

\begin{enumerate}
    \item \textbf{Synthetic Data:} Generate a dataset representing a "mystery physical signal" $f(x) = \sin(x) + 0.5 x$ on the interval $x \in [0, 10]$. Select 15 random points and add Gaussian noise with $\sigma = 0.3$.
    \item \textbf{The Kernel:} Perform a GP regression using a **Squared Exponential (RBF) kernel**:
    $$ k(x, x') = \sigma_f^2 \exp\left( -\frac{(x-x')^2}{2\ell^2} \right) $$
    You can use \texttt{scikit-learn} or implement the covariance matrix operations directly.
    \item \textbf{Uncertainty Quantification:} Plot the interpolated mean function and the $95\%$ confidence region (shaded band).
    \item \textbf{Extrapolation:} Extend the prediction domain to $x \in [0, 15]$.
    \item \textit{Discussion:} Compare the behavior of the GP error bars in the extrapolation region ($x > 10$) to what you would get from a high-order polynomial fit. Which method "knows" that it is ignorant?
\end{enumerate}


